Learned post-hoc correction of frozen flow-matching Vision-Language-Action (VLA) models is harmful by default. The problem is not the correction module itself but its unconditional application: on most control steps the VLA already produces confident, low-variance actions, and corrections applied here act as noise---degrading all four LIBERO manipulation suites on SmolVLA (0.4B) by up to $-$14.8\,pp under a K-matched controlled comparison. We propose calibrated uncertainty gating: $K{=}2$ stochastic action samples are drawn in a single batched forward pass at zero extra latency (${\sim}$170\,ms, matching $K{=}1$), and their pairwise $L_2$ distance gates a lightweight (${\sim}$327K parameter) correction module so that only genuinely uncertain steps receive learned residual correction. Selective gating recovers up to $+$13.5\,pp of the degradation caused by unconditional correction, restoring near-baseline or above-baseline performance across all four suites. The trigger rate (fraction of corrected steps) shows a positive trend with net improvement over VLA-only, providing a built-in deployment diagnostic. Experiments on $\pi_{0.5}$ (3.6B) with $n{=}3$ seeds per condition confirm that the degradation pattern persists and amplifies ($-$17.5\,pp on Long), with partial recovery from gating. The key finding is that knowing when \emph{not} to correct matters more than the correction itself.
